
\subsection*{Weak PDE flux of the actual moving remote Hartree potential}

\noindent
The absolute density-velocity estimate in the preceding supplement
uses a high-regularity norm of $\Delta z+\mathcal N_H(z)$.
For the \emph{localized remote potential itself}, this second
spatial derivative can instead be removed using the original
Hartree PDE \emph{before} taking absolute values.

Let $iz_t+\Delta z+\mathcal N_H(z)=0$,
$\mathcal N_H(z)=\kappa_H(|x|^{-2}*|z|^4)|z|^2z$ with real
$\kappa_H$, and
$Z(s,x)=\sqrt{\ell(s)}z(t(s),\ell(s)x)$,
$\ell_s=-b\ell$, $t_s=\ell^2$. The genuine critical scaling gives
\[
 Z_s=i(\Delta Z+\mathcal N_H(Z))
       -b\Lambda Z,\qquad
 \Lambda Z=\tfrac12Z+x\cdot\nabla Z.
\]
For a real, possibly moving smooth reference $V(s,x)$ set
$U=V+Z$ and $\rho=|U|^4-|V|^4$.
Fix $L(s)>0$, $R(s)=\theta L(s)$ with
$3/4\le\theta\le1$. Let $\chi\in C_c^1([0,\infty))$,
$\operatorname{supp}\chi\subset[0,1/3]$, $|\chi|\le1$ and
$C_\chi=\|\chi'\|_\infty$. The \emph{exact} separated radial
three-dimensional kernel is
\[
 K_R(q)=\frac{1}{2Rq}\log\frac{R+q}{R-q}
 =\frac1{R^2}\sum_{j\ge0}\frac{(q/R)^{2j}}{2j+1},
 \qquad K_R(0)=R^{-2}.
\]
With $w(s,q)=\chi(q/L(s))K_{R(s)}(q)$, define
\[
 Q_\chi(s,R(s))=
 4\pi\int_0^\infty q^2w(s,q)\rho(s,q)\,dq
 =\int_{\mathbb R^3}w(s,|x|)\rho(s,x)\,dx.
\]
The separated geometry $q\le L/3$, $R\ge3L/4$
and $q/R\le4/9$ yields the explicit budgets
\[
 W_0=\frac{144}{65L^2},\qquad
 W_1=\frac{256+432C_\chi}{195L^3},\qquad
 W_R=\frac{384}{65L^3}
\]
for $|w|$, $|\partial_qw|$ and
$|\partial_Rw|$, respectively.
There is no singular left Kummer-radius factor.

Writing $\Phi=|U|^2\overline U$, use the genuine PDE and
integrate the radial Laplacian by parts in $4\pi q^2dq$:
\begin{align*}
 \frac d{ds}Q_\chi(s,R(s))
 ={}&
 \int_{\mathbb R^3}
   \left(\left.\partial_sw\right|_R+
     R'\partial_Rw\right)\rho\,dx\\
 &-4\Ree\int_{\mathbb R^3}
      i\nabla(w\Phi)\cdot\nabla Z\,dx
  +4\Ree\int_{\mathbb R^3}
      iw\Phi\,\mathcal N_H(Z)\,dx\\
 &-4b\Ree\int_{\mathbb R^3}w\Phi\Lambda Z\,dx
  +4\Ree\int_{\mathbb R^3}
       w(\Phi-|V|^2\overline V)V_s\,dx .
\end{align*}
The literal cutoff derivative is
$\left.\partial_sw\right|_R
 =-(L'/L)(q/L)\chi'(q/L)K_R(q)$.
For smooth radial fields the center boundary term vanishes;
the compactly supported cutoff eliminates the outer boundary.
The genuine nonlocal Hartree source is retained, not dropped.

Put $\Omega_L=B_{L/3}(0)$, $A_L=|\Omega_L|=4\pi L^3/81$,
$B_U=\sup_{\Omega_L}(|V|+|Z|)$,
$B_V=\sup_{\Omega_L}|V|$,
$e_0=\|Z\|_{L^2(\Omega_L)}$,
$e_1=\|\nabla Z\|_{L^2(\Omega_L)}$,
$v_1=\|\nabla V\|_{L^2(\Omega_L)}$,
$v_s=\|V_s\|_{L^2(\Omega_L)}$ and
$J_\rho=\int_{\Omega_L}|\rho|\,dx$.
Since $|\Phi|\le B_U^3$ and
$|\nabla\Phi|\le3B_U^2(|\nabla V|+|\nabla Z|)$,
\begin{align*}
 \left|\frac d{ds}Q_\chi(s,R(s))\right|
 \le{}&
 \left(\frac{C_\chi|L'|}{3L}W_0+
          |R'|W_R\right)J_\rho\\
 &+4W_1B_U^3A_L^{1/2}e_1
  +12W_0B_U^2(v_1+e_1)e_1\\
 &+4W_0B_U^3\|\mathcal N_H(Z)\|_{L^1(\Omega_L)}\\
 &+4|b|W_0B_U^3A_L^{1/2}
       \left(\frac{e_0}{2}+\frac L3e_1\right)\\
 &+4W_0(B_U^3+B_V^3)A_L^{1/2}v_s,
 \qquad
 J_\rho\le4B_U^3A_L^{1/2}e_0 .
\end{align*}
The \emph{original} Hartree term is controlled using the
three-dimensional Hardy--Littlewood--Sobolev inequality:
\[
 \|\mathcal N_H(Z)\|_{L^1(\Omega_L)}
 \le |\kappa_H|C_{\mathrm{HLS}}
          A_L^{1/6}\|Z\|_{L^6(\mathbb R^3)}^7.
\]
This avoids the second-derivative radiation quantity
$\mathfrak M(\Delta z+\mathcal N_H(z))^{1/4}$ and the
absolute-density-velocity bound $J_1$ in the localized
remote channel. Unlike a local cubic replacement, it retains
the genuine nonlocal quartic potential.

\medskip
\noindent
\textbf{Verification boundary.}
The calculation is a paper-level weak-flux lemma requiring
a classical Hartree orbit (or justified weak approximation).
It does \emph{not} prove that the actual norms, moving reference,
matching radius, Kummer basis, unique real selected speeds or
the full untruncated near-diagonal and four-region Hartree
residual have the necessary temporal gain or parameter Lipschitz
control. The finite-layer PDE recursion and dynamically derived
Type--II exponent remain unproved. No new Lean CI has accepted
these estimates; original revised theorem environments remain
$11/13$. The original quartic phase cancellation is already
formalized in \texttt{ActualMomentPDEFlux.lean}, but is not a
formal audit of the present differentiated potential.
